% 14-end-to-end-trace.tex — 端到端编码追踪
% Source: chapters/14-end-to-end-annotated-encoding-trace.md

\chapter{端到端编码追踪}\label{ch:e2e-trace}

\section{本章目标}

\begin{itemize}
  \item 通过一个完整的编码示例，串联所有编码步骤
  \item 理解数据在各模块之间的流转
  \item 掌握从原始图像到码流的完整路径
  \item 理解解码重建的完整逆路径
\end{itemize}

\section{背景与标准语义}

本章以一个最小但完整的编码示例，追踪数据从输入图像到最终码流的每一步变换。选择 $2\times2$ RGB 图像配合 $NL_x=NL_y=1$，使每个子带只有一个系数，便于手工验证。

\section{追踪配置}

\textbf{输入图像}：$2\times2$ RGB（3 分量，无下采样，8-bit）

\begin{lstlisting}[language=C,numbers=none]
R 分量:     G 分量:     B 分量:
[120  80     [200 180    [ 60  40
 100  60]     160 140]     50  30]
\end{lstlisting}

像素位置 $(0,0)$：$R=120, G=200, B=60$。像素 $(1,0)$：$R=80, G=180, B=40$。等等。

\textbf{编码配置}：

\begin{table}[H]
\centering
\caption{端到端追踪编码配置}
\label{tab:e2e-config}
\begin{tabular}{@{}lll@{}}
\toprule
参数 & 值 & 说明 \\
\midrule
$NL_x$ & 1 & 水平分解 1 层 \\
$NL_y$ & 1 & 垂直分解 1 层 \\
$F_q$ & 4 & 小数位 4 bit \\
$B_w$ & 12 & 工作位宽 \\
$N_g$ & 4 & Group 大小 \\
$C_{pih}$ & 1 & RCT 颜色变换 \\
$Q_{pih}$ & 0 & Dead-zone 量化 \\
$F_s$ & 0 & 联合符号 \\
$R_m$ & 0 & ZRF（零残差标志） \\
Gain & $[2, 1, 1, 0]$ & 每分量 band 增益 \\
Priority & $[0, 1, 2, 3]$ & 每分量 band 优先级 \\
\bottomrule
\end{tabular}
\end{table}

\section{编码追踪}

\subsection{Step 1：NLT 正变换}

配置 $T_{nlt} = \mathrm{NONE}$，跳过。输入像素值保持不变。

\subsection{Step 2：MCT 正变换（RCT）}

对每个像素位置应用 RCT（\texttt{mct.c}，公式见第 5 章）：
\begin{equation}
Y = \left\lfloor \frac{R + 2G + B}{4} \right\rfloor, \quad U = B - G, \quad V = R - G
\end{equation}

\begin{table}[H]
\centering
\caption{RCT 结果}
\label{tab:e2e-rct}
\footnotesize
\setlength{\tabcolsep}{4pt}
\begin{tabular}{@{}lcccccc@{}}
\toprule
像素 & R & G & B & \multicolumn{1}{c}{Y} & \multicolumn{1}{c}{U} & \multicolumn{1}{c}{V} \\
\midrule
(0,0) & 120 & 200 & 60  & $\lfloor(120{+}400{+}60)/4\rfloor{=}\mathbf{145}$ & $60{-}200{=}\mathbf{-140}$ & $120{-}200{=}\mathbf{-80}$ \\
(1,0) & 80  & 180 & 40  & $\lfloor(80{+}360{+}40)/4\rfloor{=}\mathbf{120}$  & $40{-}180{=}\mathbf{-140}$ & $80{-}180{=}\mathbf{-100}$ \\
(0,1) & 100 & 160 & 50  & $\lfloor(100{+}320{+}50)/4\rfloor{=}\mathbf{117}$  & $50{-}160{=}\mathbf{-110}$ & $100{-}160{=}\mathbf{-60}$ \\
(1,1) & 60  & 140 & 30  & $\lfloor(60{+}280{+}30)/4\rfloor{=}\mathbf{92}$   & $30{-}140{=}\mathbf{-110}$ & $60{-}140{=}\mathbf{-80}$ \\
\bottomrule
\end{tabular}
\end{table}

RCT 后 3 个分量各 $2\times2$：

\begin{lstlisting}[numbers=none]
Y 分量:          U 分量:           V 分量:
[145  120        [-140 -140        [-80  -100
 117   92]        -110 -110]        -60   -80]
\end{lstlisting}

\subsection{Step 3：DWT 正变换（$NL_x=1, NL_y=1$）}

对每个分量做 2D DWT。由于图像是 $2\times2$，DWT 产生 4 个 $1\times1$ 子带（每个分量）。

2D 分解顺序（\texttt{dwt.c:161}）：先垂直，后水平。

\subsubsection{Y 分量（$2\times2$）}

原始矩阵（内存中行优先）：$[145, 120, 117, 92]$

\textbf{Phase 1 垂直}（对每列，$\mathit{inc} = W = 2$）：

列 0：$[145, 117]$
\begin{itemize}
  \item 奇数（边界）：$d[1] = 117 - 145 = -28$
  \item 偶数（首偶）：$s[0] = 145 + ((-28 + 1) \gg 1) = 145 + (-27 \gg 1)$
\end{itemize}

在 C 中，$-27 \gg 1 = -14$（算术右移）。$s[0] = 145 + (-14) = 131$

\begin{tcolorbox}[colback=yellow!5,colframe=yellow!60!black,title=注意]
这里涉及负数右移。C 标准规定有符号数右移是实现定义的，但主流编译器使用算术右移。libjxs 依赖此行为。
\end{tcolorbox}

列 0 结果：$[131, -28]$

列 1：$[120, 92]$
\begin{itemize}
  \item 奇数（边界）：$d[1] = 92 - 120 = -28$
  \item 偶数（首偶）：$s[0] = 120 + ((-28 + 1) \gg 1) = 120 + (-14) = 106$
\end{itemize}
列 1 结果：$[106, -28]$

垂直分解后：$[131, 106, -28, -28]$（低频行在前，高频行在后）

\textbf{Phase 1 水平}（对每行，$\mathit{inc} = 1$）：

行 0（低频行）：$[131, 106]$
\begin{itemize}
  \item 奇数（边界）：$d[1] = 106 - 131 = -25$
  \item 偶数（首偶）：$s[0] = 131 + ((-25 + 1) \gg 1) = 131 + (-12) = 119$
\end{itemize}
行 0 结果：$[119, -25]$

行 1（高频行）：$[-28, -28]$
\begin{itemize}
  \item 奇数（边界）：$d[1] = -28 - (-28) = 0$
  \item 偶数（首偶）：$s[0] = -28 + ((0 + 1) \gg 1) = -28 + 0 = -28$
\end{itemize}
行 1 结果：$[-28, 0]$

\textbf{Y 分量 DWT 结果}（4 个 $1\times1$ 子带）：

\begin{lstlisting}[numbers=none]
LL: [119]    HL: [-25]
LH: [-28]    HH: [0]
\end{lstlisting}

在内存中：$[119, -25, -28, 0]$（位置 0=LL, 1=HL, 2=LH, 3=HH）

\subsubsection{U 分量（$2\times2$）}

原始：$[-140, -140, -110, -110]$

\textbf{垂直}：
\begin{itemize}
  \item 列 0：$[-140, -110]$：$d = -110-(-140) = 30$，$s = -140 + ((30+1)\gg1) = -140+15 = \mathbf{-125}$
  \item 列 1：$[-140, -110]$：$d = 30$，$s = \mathbf{-125}$
\end{itemize}
垂直后：$[-125, -125, 30, 30]$

\textbf{水平}：
\begin{itemize}
  \item 行 0：$[-125, -125]$：$d = -125-(-125) = 0$，$s = -125+0 = \mathbf{-125}$
  \item 行 1：$[30, 30]$：$d = 30-30 = 0$，$s = 30+0 = \mathbf{30}$
\end{itemize}

\textbf{U 分量 DWT 结果}：

\begin{lstlisting}[numbers=none]
LL: [-125]    HL: [0]
LH: [30]      HH: [0]
\end{lstlisting}

\subsubsection{V 分量（$2\times2$）}

原始：$[-80, -100, -60, -80]$

\textbf{垂直}：
\begin{itemize}
  \item 列 0：$[-80, -60]$：$d = 20$，$s = -80 + 10 = \mathbf{-70}$
  \item 列 1：$[-100, -80]$：$d = 20$，$s = -100+10 = \mathbf{-90}$
\end{itemize}
垂直后：$[-70, -90, 20, 20]$

\textbf{水平}：
\begin{itemize}
  \item 行 0：$[-70, -90]$：$d = -20$，$s = -70 + ((-20+1)\gg1) = -70+(-10) = \mathbf{-80}$
  \item 行 1：$[20, 20]$：$d = 0$，$s = \mathbf{20}$
\end{itemize}

\textbf{V 分量 DWT 结果}：

\begin{lstlisting}[numbers=none]
LL: [-80]    HL: [-20]
LH: [20]     HH: [0]
\end{lstlisting}

\subsection{Step 4：IDS 构建}

\texttt{ids\_construct()} 构建空间索引。对于 $2\times2$ 输入、$NL_x=NL_y=1$：

\begin{lstlisting}[numbers=none]
每个分量: 4 bands (LL, HL, LH, HH)
每个 band: 1×1（1 个系数）
总 bands: 3 × 4 = 12
precinct 数: npx × npy = 1 × 1 = 1（整个图像只有 1 个 precinct）
\end{lstlisting}

\subsection{Step 5：Precinct 提取与 $F_q$ 转换}

\texttt{precinct\_from\_image()} 将 DWT 系数转换为 sign-magnitude 表示。

转换公式（\texttt{precinct.c}，$F_q = 4$）：
\begin{lstlisting}[language=C,numbers=none]
Fq_r = (1 << Fq) >> 1;  // = 8
if (val >= 0)
    sm = (val + Fq_r) >> Fq;
else
    sm = ((-val + Fq_r) >> Fq) | SIGN_BIT_MASK;
\end{lstlisting}

\begin{table}[H]
\centering
\caption{$F_q$ 转换结果}
\label{tab:e2e-fq}
\begin{tabular}{@{}llrll@{}}
\toprule
Band & 分量 & DWT 系数 & 计算过程 & Sign-Magnitude \\
\midrule
LL & Y & 119  & $(119+8)\gg4 = 127\gg4 = \mathbf{7}$   & 7 \\
HL & Y & $-25$ & $(25+8)\gg4 = 33\gg4 = \mathbf{2}$，负 $\to 2\,|\,0x80000000$ & 0x80000002 \\
LH & Y & $-28$ & $(28+8)\gg4 = 36\gg4 = \mathbf{2}$，负 $\to 2\,|\,S$ & 0x80000002 \\
HH & Y & 0    & $(0+8)\gg4 = 0$ & \textbf{0} \\
LL & U & $-125$ & $(125+8)\gg4 = 133\gg4 = \mathbf{8}$，负 $\to 8\,|\,S$ & 0x80000008 \\
HL & U & 0    & \textbf{0} & 0 \\
LH & U & 30   & $(30+8)\gg4 = 38\gg4 = \mathbf{2}$ & 2 \\
HH & U & 0    & \textbf{0} & 0 \\
LL & V & $-80$ & $(80+8)\gg4 = 88\gg4 = \mathbf{5}$，负 $\to 5\,|\,S$ & 0x80000005 \\
HL & V & $-20$ & $(20+8)\gg4 = 28\gg4 = \mathbf{1}$，负 $\to 1\,|\,S$ & 0x80000001 \\
LH & V & 20   & $(20+8)\gg4 = 28\gg4 = \mathbf{1}$ & 1 \\
HH & V & 0    & \textbf{0} & 0 \\
\bottomrule
\end{tabular}
\end{table}

\subsection{Step 6：GCLI 计算}

\texttt{update\_gclis()} 对每个 band 行按 $N_g=4$ 分组计算 GCLI。由于每个 band 只有 1 个系数（1 组），GCLI 直接由该系数的幅度决定：
\begin{equation}
\mathit{GCLI} = \mathrm{BSR}(\text{幅度}) + 1
\end{equation}

\begin{table}[H]
\centering
\caption{GCLI 计算结果}
\label{tab:e2e-gcli}
\begin{tabular}{@{}lrlrl@{}}
\toprule
Band & 幅度 & Binary & BSR & GCLI \\
\midrule
LL (Y) & 7  & 0b0111 & 2 & \textbf{3} \\
HL (Y) & 2  & 0b0010 & 1 & \textbf{2} \\
LH (Y) & 2  & 0b0010 & 1 & \textbf{2} \\
HH (Y) & 0  & ---    & --- & \textbf{0} \\
LL (U) & 8  & 0b1000 & 3 & \textbf{4} \\
HL (U) & 0  & ---    & --- & \textbf{0} \\
LH (U) & 2  & 0b0010 & 1 & \textbf{2} \\
HH (U) & 0  & ---    & --- & \textbf{0} \\
LL (V) & 5  & 0b0101 & 2 & \textbf{3} \\
HL (V) & 1  & 0b0001 & 0 & \textbf{1} \\
LH (V) & 1  & 0b0001 & 0 & \textbf{1} \\
HH (V) & 0  & ---    & --- & \textbf{0} \\
\bottomrule
\end{tabular}
\end{table}

GCLI 表（12 个 band 的 GCLI 值）：$[3, 2, 2, 0, 4, 0, 2, 0, 3, 1, 1, 0]$

\subsection{Step 7：码率控制}

假设总预算 200 bit，图像只有 1 个 precinct，全部分配给当前 precinct。

$\mathit{GTLI} = \mathrm{clip}(Q - \mathit{gain} - [\mathit{priority} < R],\; 0,\; 15)$

编码量计算（简化）：$Q=0, R=0$ 时 $\mathit{GTLI} = [0, 0, 0, 0]$，编码量 $\approx 120$ bit $< 200$，直接进入阶段 2。

最终：$Q=0, R=\max$（无量化）。

\subsection{Step 8：量化}

$Q=0, R=\max$ 时 $\mathit{GTLI}=[0, 0, 0, 0]$，不做截断。所有系数保持原值。

\subsection{Step 9：打包}

以 LL(Y) 为例（$\mathit{GCLI}=3$，幅度=7，正数）：
\begin{itemize}
  \item 符号位：0（正）
  \item Bitplane 2（$\mathit{GCLI}-1$）：7 的 bit 2 = 1
  \item Bitplane 1：7 的 bit 1 = 1
  \item Bitplane 0（$\mathit{GTLI}$）：7 的 bit 0 = 1
  \item 编码：\texttt{0 1 1 1} = 4 bit
\end{itemize}

\subsection{Step 10：写入码流}

最终码流结构：

\begin{lstlisting}[numbers=none]
SOC (0xFF10)
PIH (0xFF12) — 图像参数
CDT (0xFF13) — 分量描述: [8,1,1] × 3
WGT (0xFF14) — band gains/priorities
SLH (0xFF20) — slice header
Precinct data — 编码后的系数
EOC (0xFF11)
\end{lstlisting}

\section{解码追踪（逆路径）}

\subsection{码流解析}

\begin{lstlisting}[numbers=none]
xs_parse_head() → 提取 PIH, CDT, WGT
ids_construct() → 重建空间索引
\end{lstlisting}

\subsection{Precinct 解包}

\begin{lstlisting}[numbers=none]
unpack_precinct():
  1. 读取 precinct header → Q=0, R=max, GTLI=[0,0,0,0]
  2. 解码 subpacket:
     - GCLI: [3, 2, 2, 0, 4, 0, 2, 0, 3, 1, 1, 0]
     - 系数: sign-magnitude 值
\end{lstlisting}

\subsection{逆量化}

$\mathit{GTLI}=0$，无截断。sign-magnitude 值直接使用。

$\mathit{LL(Y)}$: sm=7（正）$\to$ DWT 系数 $= 7 \ll 4 = \mathbf{112}$

$\mathit{HL(Y)}$: sm=0x80000002（负）$\to$ DWT 系数 $= -(2 \ll 4) = \mathbf{-32}$

\begin{tcolorbox}[colback=yellow!5,colframe=yellow!60!black,title=精度损失]
原始 DWT 系数为 119 和 $-25$，重建为 112 和 $-32$。这是因为 $F_q = 4$ 的定点化引入了 $\pm 8$ 以内的误差。
\end{tcolorbox}

\subsection{逆 DWT}

对 Y 分量做逆 DWT。DWT 结果：$[112, -32, -32, 0]$（LL, HL, LH, HH）。

逆变换（\texttt{dwt.c:143}）：先逆水平，后逆垂直。

\textbf{逆水平}（对每行）：

行 0：$[112, -32]$
\begin{itemize}
  \item 恢复偶数：$x[0] = 112 - ((-32+1)\gg1) = 112 - (-16) = \mathbf{128}$
  \item 恢复奇数：$x[1] = -32 + 128 = \mathbf{96}$
\end{itemize}

行 1：$[-32, 0]$
\begin{itemize}
  \item 恢复偶数：$x[0] = -32 - ((0+1)\gg1) = -32 - 0 = \mathbf{-32}$
  \item 恢复奇数：$x[1] = 0 + (-32) = \mathbf{-32}$
\end{itemize}

逆水平后：$[128, 96, -32, -32]$

\textbf{逆垂直}（对每列）：

列 0：$[128, -32]$
\begin{itemize}
  \item 恢复偶数：$x[0] = 128 - ((-32+1)\gg1) = 128 - (-16) = \mathbf{144}$
  \item 恢复奇数：$x[1] = -32 + 144 = \mathbf{112}$
\end{itemize}

列 1：$[96, -32]$
\begin{itemize}
  \item 恢复偶数：$x[0] = 96 - ((-32+1)\gg1) = 96 - (-16) = \mathbf{112}$
  \item 恢复奇数：$x[1] = -32 + 112 = \mathbf{80}$
\end{itemize}

\textbf{Y 分量恢复}：$[144, 112, 112, 80]$

原始 Y 值：$[145, 120, 117, 92]$。恢复值：$[144, 112, 112, 80]$。

误差：$-1, -8, -5, -12$。误差来源是 $F_q = 4$ 的定点化截断。

\subsection{逆 MCT（RCT）}

对每个像素位置：$G = Y - \lfloor(U + V)/4\rfloor,\; R = V + G,\; B = U + G$

像素 $(0,0)$：$Y=144, U \approx -136, V \approx -80$
\begin{align}
G &= 144 - ((-136 + (-80)) \gg 2) = 144 - (-54) = 198 \\
R &= -80 + 198 = 118 \\
B &= -136 + 198 = 62
\end{align}

原始值：$R=120, G=200, B=60$。重建值：$R=118, G=198, B=62$。

误差：$\Delta R = -2, \Delta G = -2, \Delta B = +2$。在 8-bit 可视范围内，误差小于 $3/255$，不可感知。

\subsection{逆 NLT}

$T_{nlt} = \mathrm{NONE}$，跳过。

\section{数据流转图}

\begin{figure}[H]
\centering
\begin{tikzpicture}[
    node distance=0.7cm and 0.6cm,
    block/.style={rectangle, draw=structurecolor!60, fill=structurecolor!10,
                  text width=1.5cm, align=center, minimum height=0.65cm,
                  font=\footnotesize, inner sep=3pt},
    arrow/.style={-{Stealth[length=2mm]}, thick, draw=structurecolor!60}
]
    % Row 1
    \node[block] (rgb) {RGB};
    \node[block, right=of rgb] (nlt) {NLT};
    \node[block, right=of nlt] (mct) {MCT};
    \node[block, right=of mct] (dwt) {DWT};
    \node[block, right=of dwt] (prec) {Precinct};

    % Row 2
    \node[block, below=0.9cm of rgb] (gcli) {GCLI};
    \node[block, right=of gcli] (rc) {RC};
    \node[block, right=of rc] (quant) {Quant};
    \node[block, right=of quant] (pack) {Pack};
    \node[block, right=of pack] (cs) {Code-stream};

    % Row 1 arrows
    \draw[arrow] (rgb) -- (nlt);
    \draw[arrow] (nlt) -- (mct);
    \draw[arrow] (mct) -- (dwt);
    \draw[arrow] (dwt) -- (prec);

    % Row 2 arrows
    \draw[arrow] (gcli) -- (rc);
    \draw[arrow] (rc) -- (quant);
    \draw[arrow] (quant) -- (pack);
    \draw[arrow] (pack) -- (cs);

    % Inter-row arrow
    \draw[arrow] (prec.south) -- ++(0,-0.3) -| (gcli.north);

    % Label
    \node[above=0.1cm of rgb, font=\footnotesize\bfseries] {编码路径};
\end{tikzpicture}
\caption{端到端编码数据流}
\label{fig:e2e-flow}
\end{figure}

\section{精度分析}

本例中 $F_q = 4$，即 DWT 系数以 16 为单位量化。每一步引入的误差：

\begin{table}[H]
\centering
\caption{精度分析}
\label{tab:e2e-precision}
\begin{tabular}{@{}lll@{}}
\toprule
步骤 & 误差来源 & 典型幅度 \\
\midrule
MCT (RCT) & 除以 4 的截断 & $\pm 1$ \\
DWT & 整数 lifting 的右移 & $\pm 1$ per step \\
$F_q$ 转换 & 右移 4 位 & $\pm 8$ (DWT 域) \\
量化 (GTLI) & bitplane 截断 & $2^{\mathit{GTLI}}$ per coefficient \\
逆过程 & 累积误差 & 最终 $\pm 3$ per pixel \\
\bottomrule
\end{tabular}
\end{table}

在本例中（$Q=0$，无量化），主要误差来自 $F_q$ 转换和 DWT 整数运算。

\section{本章小结}

通过这个 $2\times2$ RGB 的端到端追踪，我们展示了 JPEG XS 编码的完整流程：NLT $\to$ MCT(RCT) $\to$ DWT $\to$ precinct 提取 $\to$ GCLI 计算 $\to$ 码率控制 $\to$ 量化 $\to$ 打包。每一步的中间结果都可以通过手工计算验证。解码路径严格逆向执行，恢复出近似的图像。

关键发现：
\begin{itemize}
  \item RCT 在像素域做颜色去相关，YUV 分量的动态范围比 RGB 更大
  \item $2\times2$ 图像的 DWT 退化为 4 个 $1\times1$ 子带，但 lifting 公式仍然完整执行
  \item $F_q = 4$ 的定点化引入 $\pm 8$ 的 DWT 域误差，逆变换后约为 $\pm 3$ 像素误差
  \item 无量化（$Q=0$）时，主要误差来源是整数 lifting 和 $F_q$ 截断
\end{itemize}
